\section*{Pompa Howland}

\noindent\colorbox{orange}{
  \begin{minipage}{1.0\linewidth}
    \subsection*{Richiesta (a)}
    Nell'ipotesi che l'operazionale sia ideale e la sua uscita lineare ed assunti i valori delle
    resistenze pari a quelli nominali (e quindi $R_1=...=R_4=10~\text{kW}$), determinare l'impedenza
    di ingresso del blocco tratteggiato in figura e verificare l'accordo con quanto ottenuto al
    punto 1c.
  \end{minipage}
}
\subsection*{Risposta (a)}
Si studia la stabilit\'a del circuito
\begin{equation*}
  \begin{gathered}
    V_B^{(n+1)} = \epsilon A_d(V_A - V_B^{(n)}/2) + (1 - \epsilon)V_B^{(n)} \\ \implies
    V_B^{(n+1)} = \epsilon A_dV_A + \boxed{[1 - \epsilon(1 + A_d/2)]}V_B^{(n)}
  \end{gathered}
\end{equation*}
siccome $\forall A_d \exists \epsilon > 0 \colon 1 - \epsilon(1 + A_d/2) < 1$ allora
il circuito \'e \textit{lineare/stabile} e dunque $V_B \equiv V_B^{(n+1)} = V_B^{(n)}$
converger\'a ad un qualche valore finito.
\begin{equation*}
  \begin{gathered}
    V_B = \epsilon A_dV_A + [1 - \epsilon(1 + A_d/2)]V_B \\ \implies
    V_B = \frac{A_dV_A}{(1 + A_d/2)}
  \end{gathered}
\end{equation*}
nell'ipotesi di idealit\'a dell'op-amp in termini di guadagno differenziale si ottiene
\begin{equation*}
  \boxed{V_B = 2V_A}
\end{equation*}
dunque si calcola la corrente che scorre attraverso $R_2$ (in funzione di $V_A$) ed infine
l'impedenza equivalente
\begin{equation*}
  I_A = \frac{V_A - V_B}{R} \implies \colorbox{blu}{\boxed{Z_{eq} \equiv \frac{V_A}{I_A} = -R}}
\end{equation*}

\noindent\colorbox{orange}{
  \begin{minipage}{1.0\linewidth}
    \subsection*{Richiesta (b)}
    Ottenere un'espressione per $V_A$ e $V_B$ ed i valori attesi per $V_S=1~\text{V}$ ed $R_L=10~\text{k}\Omega$,
    confrontandoli con le misure di cui al punto 1b.
  \end{minipage}
}
\subsection*{Risposta (b)}
Si calcola $V_A$ sostituendo il blocco con un resistore equivalente con \textit{resistenza negativa} da cui
si ricava l'espressione in funzione di $V_S$ e $R_L$ attraverso la KCL
\begin{equation*}
  \begin{gathered}
    \frac{V_S - V_A}{R} = \frac{V_A}{R_L} + \frac{V_A}{-R} \\ \implies
    \colorbox{blu}{\boxed{V_A = \frac{R_L}{R}V_S}}
  \end{gathered}
\end{equation*}
la relazione fra $V_A$ e $V_B$ \'e nota da prima, dunque
\begin{equation*}
  \colorbox{blu}{\boxed{V_B = 2\frac{R_L}{R}V_S}}
\end{equation*}
i valori attesi di conseguenza
\begin{equation*}
  \colorbox{blu}{\boxed{
    \begin{gathered}
      V_A = 1~\text{V} \\
      V_B = 2~\text{V}
    \end{gathered}
  }}
\end{equation*}

\noindent\colorbox{orange}{
  \begin{minipage}{1.0\linewidth}
    \subsection*{Richiesta (c)}
    Dalle precedenti espressioni determinare un valore atteso per la transconduttanza
    dell'amplificatore in termini dei componenti circuitali, verificando che sia compatibile con
    quanto ottenuto al punto 1e ed indipendente dalla resistenza di carico come mostrato
    sperimentalmente al punto 1f.
  \end{minipage}
}
\subsection*{Risposta (c)}
La definizione di \textit{transconduttanza} riportata nella scheda completa dell'esercizio \'e la
seguente
\begin{equation*}
  G_m = \frac{I_{OUT}}{V_S}
\end{equation*}
la corrente \'e banale da calcolare siccome si ha un'espressione di $V_A$ e di conseguenza anche
la transconduttanza
\begin{equation*}
  I_{OUT} = \frac{V_A}{R_L} \implies \colorbox{blu}{\boxed{G_m = \frac{1}{R}}}
\end{equation*}

\noindent\colorbox{orange}{
  \begin{minipage}{1.0\linewidth}
    \subsection*{Richiesta (d)}
    Nelle stesse ipotesi di cui al punto 2a determinare, a partire dalla sua definizione, l'impedenza
    di uscita del circuito (N.B. i terminali di uscita devono essere identificati con quelli del
    carico).
  \end{minipage}
}
\subsection*{Risposta (d)}
\begin{equation*}
  Z_{out} = \left(\frac{1}{R_1} + \frac{1}{-R_2}\right)^{-1} = \frac{R_1R_2}{R_1 - R_2}
\end{equation*}

\noindent\colorbox{orange}{
  \begin{minipage}{1.0\linewidth}
    \subsection*{Richiesta (e)}
    Discutere come il risultato ottenuto al punto precedente sia correlato con l'indipendenza di
    $I_{OUT}$ da $R_L$.
  \end{minipage}
}
\subsection*{Risposta (e)}
